% GENERATED FILE. Edit canonical lessons, then run npm run generate:books.
% canonical-source-hash: fnv1a64:6888fb911d5d6fd7
% canonical-lessons: FR-C366-benevole, FR-C366-benevolat, FR-C366-association, FR-C366-solidarite, FR-C366-communaute

\chapter{Donner un coup de main}
\label{ch:fr-donner-un-coup-de-main}
\hlchaptermodality{366}

\begin{chapteropening}
\textbf{By the end of this chapter:} \emph{I can talk about helping others: volunteering, a volunteer, an association, solidarity and the community.}

Complete the last lesson of chapter 366: I can talk about helping others: volunteering, a volunteer, an association, solidarity and the community.
\end{chapteropening}

\section[le bénévole]{le bénévole — the volunteer; unpaid, voluntary}

\label{lesson:FR-C366-benevole}



\begin{quote}
Before the new one: say the French for to discriminate, then the French for the immigration.
\end{quote}

\subsection*{You'll want to know: le bénévole}
\textbf{le bénévole} — \textquotedblleft{}the volunteer; unpaid, voluntary\textquotedblright{}.

\begin{quote}
Les bénévoles servent des repas aux personnes âgées.
\end{quote}

\begin{quote}
\emph{The volunteers serve meals to elderly people.}
\end{quote}

The same form for a man or a woman. As an adjective: \emph{un travail bénévole}.

\begin{cousinweb}
\emph{bénévole} is a learned borrowing from Latin \emph{benevolus} \textquotedblleft{}well-wishing, kind\textquotedblright{} (Lewis \& Short), built from \emph{bene} \textquotedblleft{}well\textquotedblright{} and \emph{velle} \textquotedblleft{}to want\textquotedblright{}.
\end{cousinweb}

\begin{grammarlens}[title={What you've built}]
One more word for telling, explaining and describing at B1.
\end{grammarlens}

\subsection*{Your turn}
\begin{itemize}
\raggedright
  \item \emph{Say it:} \emph{le bénévole}
  \item \emph{Say it:} \emph{le bénévole}, once more
  \item \emph{Say it:} a sentence of your own with \emph{le bénévole}
  \item \emph{Recall:} say the French for to discriminate, then the French for the immigration, then say \emph{le bénévole} again
\end{itemize}

\begin{tcolorbox}[breakable,colback=teal!4,colframe=teal!35!black,title={Before you move on}]
What does le bénévole mean? (\textquotedblleft{}The volunteer; unpaid, voluntary\textquotedblright{}.) Say it once more.
\end{tcolorbox}
\section[le bénévolat]{le bénévolat — the volunteering, unpaid work}

\label{lesson:FR-C366-benevolat}



\begin{quote}
Before the new one: say the French for the immigration, then the French for the volunteer.
\end{quote}

\subsection*{You'll want to know: le bénévolat}
\textbf{le bénévolat} — \textquotedblleft{}the volunteering, unpaid work\textquotedblright{}.

\begin{quote}
Le bénévolat prend tous ses samedis.
\end{quote}

\begin{quote}
\emph{Volunteering takes up all her Saturdays.}
\end{quote}

\emph{Faire du bénévolat} — to do voluntary work.

\begin{cousinweb}
\emph{bénévolat} is built from \emph{bénévole} and \emph{-at}. \emph{bénévole} comes from Latin \emph{benevolus} \textquotedblleft{}well-wishing\textquotedblright{}, from \emph{bene} \textquotedblleft{}well\textquotedblright{} and \emph{velle} \textquotedblleft{}to want\textquotedblright{} (Lewis \& Short).
\end{cousinweb}

\begin{grammarlens}[title={What you've built}]
One more word for telling, explaining and describing at B1.
\end{grammarlens}

\subsection*{Your turn}
\begin{itemize}
\raggedright
  \item \emph{Say it:} \emph{le bénévolat}
  \item \emph{Say it:} \emph{le bénévolat}, once more
  \item \emph{Say it:} a sentence of your own with \emph{le bénévolat}
  \item \emph{Recall:} say the French for the immigration, then the French for the volunteer, then say \emph{le bénévolat} again
\end{itemize}

\begin{tcolorbox}[breakable,colback=teal!4,colframe=teal!35!black,title={Before you move on}]
What does le bénévolat mean? (\textquotedblleft{}The volunteering, unpaid work\textquotedblright{}.) Say it once more.
\end{tcolorbox}
\section[l'association]{l'association — the association, the club, the charity}

\label{lesson:FR-C366-association}



\begin{quote}
Before the new one: say the French for the volunteer, then the French for the volunteering.
\end{quote}

\subsection*{You'll want to know: l'association}
\textbf{l'association} — \textquotedblleft{}the association, the club, the charity\textquotedblright{}.

\begin{quote}
Je travaille dans une association qui aide les personnes âgées.
\end{quote}

\begin{quote}
\emph{I work for an association that helps elderly people.}
\end{quote}

In France, \emph{une association} is a non-profit group of any kind. Feminine.

\begin{cousinweb}
\emph{association} is built on \emph{associer} (EtymDB), from Latin \emph{associāre} \textquotedblleft{}to join to\textquotedblright{}, built on \emph{socius} \textquotedblleft{}companion\textquotedblright{} (Lewis \& Short).
\end{cousinweb}

\begin{grammarlens}[title={What you've built}]
One more word for telling, explaining and describing at B1.
\end{grammarlens}

\subsection*{Your turn}
\begin{itemize}
\raggedright
  \item \emph{Say it:} \emph{l'association}
  \item \emph{Say it:} \emph{l'association}, once more
  \item \emph{Say it:} a sentence of your own with \emph{l'association}
  \item \emph{Recall:} say the French for the volunteer, then the French for the volunteering, then say \emph{l'association} again
\end{itemize}

\begin{tcolorbox}[breakable,colback=teal!4,colframe=teal!35!black,title={Before you move on}]
What does l'association mean? (\textquotedblleft{}The association, the club, the charity\textquotedblright{}.) Say it once more.
\end{tcolorbox}
\section[la solidarité]{la solidarité — the solidarity}

\label{lesson:FR-C366-solidarite}



\begin{quote}
Before the new one: say the French for the volunteering, then the French for the association.
\end{quote}

\subsection*{You'll want to know: la solidarité}
\textbf{la solidarité} — \textquotedblleft{}the solidarity\textquotedblright{}.

\begin{quote}
Après l'incendie, il y a eu beaucoup de solidarité dans le village.
\end{quote}

\begin{quote}
\emph{After the fire, there was a lot of solidarity in the village.}
\end{quote}

From \emph{solidaire}, standing together.

\begin{cousinweb}
\emph{solidarité} is built on \emph{solidaire} (EtymDB), from the Latin legal phrase \emph{in solidum} \textquotedblleft{}for the whole sum\textquotedblright{}, from \emph{solidus} \textquotedblleft{}solid, whole\textquotedblright{} (Lewis \& Short).
\end{cousinweb}

\begin{grammarlens}[title={What you've built}]
One more word for telling, explaining and describing at B1.
\end{grammarlens}

\subsection*{Your turn}
\begin{itemize}
\raggedright
  \item \emph{Say it:} \emph{la solidarité}
  \item \emph{Say it:} \emph{la solidarité}, once more
  \item \emph{Say it:} a sentence of your own with \emph{la solidarité}
  \item \emph{Recall:} say the French for the volunteering, then the French for the association, then say \emph{la solidarité} again
\end{itemize}

\begin{tcolorbox}[breakable,colback=teal!4,colframe=teal!35!black,title={Before you move on}]
What does la solidarité mean? (\textquotedblleft{}The solidarity\textquotedblright{}.) Say it once more.
\end{tcolorbox}
\section[la communauté]{la communauté — the community}

\label{lesson:FR-C366-communaute}



\begin{quote}
Before the new one: say the French for the association, then the French for the solidarity.
\end{quote}

\subsection*{You'll want to know: la communauté}
\textbf{la communauté} — \textquotedblleft{}the community\textquotedblright{}.

\begin{quote}
Toute la communauté a aidé cette famille.
\end{quote}

\begin{quote}
\emph{The whole community helped this family.}
\end{quote}

\emph{La vie en communauté} — living together in a group.

\begin{cousinweb}
\emph{communauté} is usually explained as built on \emph{commun} \textquotedblleft{}common\textquotedblright{}, which is inherited from Latin \emph{commūnis} \textquotedblleft{}shared, common\textquotedblright{} (EtymDB), after Latin \emph{commūnitās} \textquotedblleft{}fellowship\textquotedblright{} (Lewis \& Short).
\end{cousinweb}

\begin{grammarlens}[title={What you've built}]
One more word for telling, explaining and describing at B1.
\end{grammarlens}

\subsection*{Your turn}
\begin{itemize}
\raggedright
  \item \emph{Say it:} \emph{la communauté}
  \item \emph{Say it:} \emph{la communauté}, once more
  \item \emph{Say it:} a sentence of your own with \emph{la communauté}
  \item \emph{Recall:} say the French for the association, then the French for the solidarity, then say \emph{la communauté} again
\end{itemize}

\begin{tcolorbox}[breakable,colback=teal!4,colframe=teal!35!black,title={Before you move on}]
What does la communauté mean? (\textquotedblleft{}The community\textquotedblright{}.) Say it once more.
\end{tcolorbox}
